\begin{cuzchapter}{系统设计与实现}{chap:design}
\usetikzlibrary{calc}

    %---------------------------------------------------------------------------%
    %->> 系统设计与实现章节写作指南
    %---------------------------------------------------------------------------%
    % 本章主要介绍APPT平台的设计思路、架构和实现方法
    % 一个完整的系统设计与实现章节应包含以下几个部分：
    % 1. 设计目标：明确系统的设计目标和原则
    % 2. 系统架构：介绍系统的整体架构和各组件的功能
    % 3. 关键技术：详细说明系统实现中使用的关键技术
    % 4. 实现细节：介绍系统的具体实现方法和过程
    % 5. 测试与评估：说明系统的测试方法和评估结果

    \section{设计目标与原则}\label{sec:design-goals}

    \subsection{设计目标}

    APPT（AI驱动的幻灯片协作编辑平台）的设计目标是创建一个功能完善、易用性强的智能幻灯片协作平台，具体目标包括：

    \begin{itemize}
        \item \textbf{实时协作支持}：基于现代Web技术实现多用户实时协作编辑，支持实时光标追踪和在线状态显示。

        \item \textbf{权限管理完善}：设计基于RBAC模型的四级权限管理系统，实现细粒度的文档访问控制。

        \item \textbf{AI辅助功能}：集成大语言模型API，支持幻灯片内容生成、语法适配和智能建议。

        \item \textbf{版本控制简化}：实现直观的版本控制系统，提供版本历史查看、差异对比和一键回滚功能。

        \item \textbf{用户体验优化}：降低技术使用门槛，为非技术用户提供友好的操作界面和简化的协作流程。

        \item \textbf{系统性能稳定}：确保系统在高并发场景下的稳定性和响应速度。
    \end{itemize}

    \subsection{设计原则}

    在系统设计过程中，我们遵循以下设计原则：

    \begin{itemize}
        \item \textbf{模块化设计}：系统划分为独立的功能模块，便于维护和扩展。

        \item \textbf{前后端分离}：实现前端展示层与后端业务层的完全分离，提高系统的灵活性和可维护性。

        \item \textbf{实时优先}：优先考虑实时协作体验，采用WebSocket协议和CRDT算法确保低延迟的协作效果。

        \item \textbf{安全可靠}：采用JWT认证和RBAC权限控制，确保用户数据和文档内容的安全性。

        \item \textbf{用户友好}：设计直观的用户界面和简化的工作流程，降低非技术用户的学习成本。

        \item \textbf{开放扩展}：提供插件化架构，支持功能模块的灵活扩展和第三方集成。
    \end{itemize}

    \section{系统总体架构}\label{sec:system-architecture}

    APPT平台采用前后端分离的架构，整体架构如图 \ref{fig:appt-architecture} 所示。系统分为前端展示层、后端业务层和数据存储层三个主要层次。

    \begin{figure}[htbp]
        \centering
        \begin{tikzpicture}[
            node distance=2cm and 3cm,
            layer/.style={rectangle, draw, rounded corners, fill=white, align=center, minimum width=4cm, minimum height=1.2cm},
            component/.style={rectangle, draw, fill=white, align=center, minimum width=3.5cm, minimum height=1.2cm},
            arrow/.style={->, >=stealth, thick},
            label/.style={text width=4cm, align=center, font=\small}
        ]

        \begin{scope}
        % 定义层节点
        \node[layer, fill=blue!10] (frontend) {前端展示层};
        \node[layer, fill=green!10, below=3.5cm of frontend] (backend) {后端业务层};
        \node[layer, fill=orange!10, below=3.5cm of backend] (storage) {数据存储层};

        % 前端组件
        \node[component, fill=blue!5, left=2.5cm of frontend] (react) {React应用\\React 19 + TypeScript + Vite};
        \node[component, fill=blue!5, right=2.5cm of frontend] (editor) {CodeMirror 6编辑器\\Markdown语法高亮};
        \node[component, fill=blue!5, above=1.5cm of react] (yjs) {Yjs客户端\\实时协作};
        \node[component, fill=blue!5, above=1.5cm of editor] (slidev) {Slidev预览组件\\实时预览};

        % 后端组件
        \node[component, fill=green!5, left=2.5cm of backend] (nest) {NestJS框架\\RESTful API};
        \node[component, fill=green!5, right=2.5cm of backend] (hocus) {Hocuspocus服务器\\WebSocket协作};
        \node[component, fill=green!5, above=1.5cm of nest] (business) {业务模块\\用户/幻灯片/权限管理};
        \node[component, fill=green!5, above=1.5cm of hocus] (auth) {认证授权\\JWT + RBAC};

        % 存储组件
        \node[component, fill=orange!5, left=2cm of storage] (mysql) {MySQL数据库\\结构化数据存储};
        \node[component, fill=orange!5, right=2cm of storage] (oss) {对象存储(OSS)\\图片与静态资源};

        % 连接线
        \draw[arrow] (react) -- (frontend);
        \draw[arrow] (editor) -- (frontend);
        \draw[arrow] (yjs) -- (frontend);
        \draw[arrow] (slidev) -- (frontend);

        \draw[arrow] (frontend) -- node[right] {HTTP/WebSocket} (backend);

        \draw[arrow] (nest) -- (backend);
        \draw[arrow] (hocus) -- (backend);
        \draw[arrow] (business) -- (backend);
        \draw[arrow] (auth) -- (backend);

        \draw[arrow] (backend) -- node[right] {数据持久化} (storage);

        \draw[arrow] (mysql) -- (storage);
        \draw[arrow] (oss) -- (storage);

        % 标签
        \node[label, above=0.5cm of frontend] {\textbf{用户界面与交互}};
        \node[label, above=0.5cm of backend] {\textbf{业务逻辑处理}};
        \node[label, above=0.5cm of storage] {\textbf{数据持久化存储}};

        \end{scope}
        \end{tikzpicture}
        \caption{APPT平台整体架构图}
        \label{fig:appt-architecture}
    \end{figure}

    \subsection{前端展示层}

    前端展示层负责用户界面呈现和交互逻辑处理，主要包括：

    \begin{itemize}
        \item \textbf{React应用}：基于React 19 + TypeScript + Vite构建的单页应用，提供响应式用户界面。

        \item \textbf{CodeMirror 6编辑器}：集成的代码编辑器，支持Markdown语法高亮和实时预览。

        \item \textbf{Yjs客户端}：实时协作客户端，通过WebSocket连接后端协作服务器。

        \item \textbf{Slidev预览组件}：实时幻灯片预览组件，支持Dev/Build双模式预览。
    \end{itemize}

    \subsection{后端业务层}

    后端业务层负责处理业务逻辑和数据持久化，主要包括：

    \begin{itemize}
        \item \textbf{NestJS框架}：基于Node.js的渐进式框架，提供RESTful API和WebSocket服务。

        \item \textbf{Hocuspocus服务器}：WebSocket协作服务器，处理Yjs文档的状态同步。

        \item \textbf{业务模块}：包括用户管理、幻灯片管理、权限控制、AI辅助、版本控制等功能模块。

        \item \textbf{认证授权}：JWT认证系统和RBAC权限控制模块。
    \end{itemize}

    \subsection{数据存储层}

    数据存储层负责数据的持久化存储，主要包括：

    \begin{itemize}
        \item \textbf{MySQL数据库}：关系型数据库，存储用户信息、幻灯片内容、权限数据等结构化数据。

        \item \textbf{对象存储}：阿里云OSS，用于存储上传的图片和其他静态资源。
    \end{itemize}

    \section{前端架构设计}\label{sec:frontend-architecture}

    \subsection{技术栈选择}

    前端采用现代Web技术栈，具体选择如下：

    \begin{itemize}
        \item \textbf{React 19.2.4}：最新React版本，提供高效的用户界面渲染和状态管理。

        \item \textbf{TypeScript 5.8}：类型安全的JavaScript超集，提高代码质量和开发效率。

        \item \textbf{Vite 6.2}：下一代前端构建工具，提供极速的开发体验和优化的生产构建。

        \item \textbf{React Router v7}：客户端路由管理，支持动态路由和嵌套路由。

        \item \textbf{CodeMirror 6}：现代化的代码编辑器，支持丰富的编辑功能和插件系统。
    \end{itemize}

    \subsection{组件化设计}

    前端采用组件化设计模式，将系统功能拆分为可复用的组件模块：

    \begin{itemize}
        \item \textbf{布局组件}：DashboardLayout、ResizablePanels等，负责页面布局和窗口管理。

        \item \textbf{编辑器组件}：EditorHeader、EditorSidebar、EditorPreview等，实现幻灯片编辑功能。

        \item \textbf{权限组件}：PermissionModal、CollaboratorModal等，管理用户权限和协作成员。

        \item \textbf{AI组件}：AIPanel、AISuggestion等，提供AI辅助功能界面。

        \item \textbf{版本控制组件}：VersionHistory、DiffViewer等，展示版本历史和内容差异。
    \end{itemize}

    \subsection{状态管理}

    前端采用分层的状态管理策略：

    \begin{itemize}
        \item \textbf{本地状态}：使用React的useState和useReducer管理组件内部状态。

        \item \textbf{全局状态}：使用Context API管理跨组件的全局状态，如用户认证状态、主题设置等。

        \item \textbf{服务器状态}：使用React Query管理从服务器获取的数据状态，支持缓存、重试和乐观更新。

        \item \textbf{协作状态}：通过Yjs管理文档的实时协作状态，包括文档内容、光标位置和用户状态。
    \end{itemize}

    \subsection{实时协作实现}

    前端实时协作功能基于Yjs库实现，具体实现方案如下：

    \begin{itemize}
        \item \textbf{文档绑定}：将CodeMirror编辑器与Yjs文档实例绑定，实现编辑操作的实时同步。

        \item \textbf{WebSocket连接}：通过@hocuspocus/provider建立与Hocuspocus服务器的WebSocket连接。

        \item \textbf{光标追踪}：使用Yjs的awareness功能实现实时光标位置显示和用户状态同步。

        \item \textbf{冲突解决}：利用CRDT算法自动解决并发编辑冲突，确保最终一致性。
    \end{itemize}

    \section{后端架构设计}\label{sec:backend-architecture}

    \subsection{NestJS框架设计}

    后端采用NestJS框架构建，具有以下特点：

    \begin{itemize}
        \item \textbf{模块化架构}：将系统功能划分为独立的模块，每个模块包含控制器、服务、实体等组件。

        \item \textbf{依赖注入}：使用依赖注入容器管理组件之间的依赖关系，提高代码的可测试性和可维护性。

        \item \textbf{中间件管道}：通过中间件管道处理请求和响应，支持身份验证、日志记录、错误处理等功能。

        \item \textbf{拦截器}：使用拦截器实现数据转换、异常处理和性能监控。
    \end{itemize}

    \subsection{RESTful API设计}

    后端提供RESTful API接口，按照资源进行组织：

    \begin{itemize}
        \item \textbf{认证接口}（/api/auth/*）：用户登录、注册、令牌刷新等认证相关操作。

        \item \textbf{用户管理}（/api/users/*）：用户信息查询、更新、密码重置等操作。

        \item \textbf{空间管理}（/api/slide-spaces/*）：幻灯片空间的创建、查询、更新、删除及成员管理。

        \item \textbf{幻灯片操作}（/api/slides/*）：幻灯片的CRUD操作、内容保存、权限设置等。

        \item \textbf{版本控制}（/api/versions/*）：版本历史查询、版本创建、内容回滚等。

        \item \textbf{评论系统}（/api/comments/*）：评论的增删改查及树形结构管理。

        \item \textbf{AI辅助}（/api/ai/*）：AI内容生成、图像生成、智能建议等功能。
    \end{itemize}

    \subsection{WebSocket协作服务}

    后端通过Hocuspocus服务器提供WebSocket协作服务：

    \begin{itemize}
        \item \textbf{连接管理}：管理多个客户端的WebSocket连接，处理连接建立、维持和断开。

        \item \textbf{文档同步}：同步Yjs文档的状态变化，确保所有客户端保持一致的文档状态。

        \item \textbf{权限验证}：在文档加载和操作时验证用户权限，确保只有授权用户才能访问和编辑文档。

        \item \textbf{数据持久化}：定期将协作文档内容保存到数据库，防止数据丢失。
    \end{itemize}

    \subsection{安全与认证}

    后端采用多层安全防护机制：

    \begin{itemize}
        \item \textbf{JWT认证}：使用JSON Web Token进行用户身份认证，支持跨域访问和单点登录。

        \item \textbf{RBAC权限控制}：基于角色的访问控制，定义Owner/Editor/Commenter/Viewer四级权限。

        \item \textbf{数据验证}：使用class-validator进行请求参数验证，防止恶意数据输入。

        \item \textbf{密码加密}：使用bcrypt算法对用户密码进行哈希存储，确保密码安全。
    \end{itemize}

    \section{数据库设计}\label{sec:database-design}

    \subsection{数据库选型}

    考虑到APPT平台的数据特点和性能需求，选择MySQL作为主要数据库：

    \begin{itemize}
        \item \textbf{关系型数据}：用户、幻灯片、权限等数据具有明确的关系结构，适合使用关系型数据库。

        \item \textbf{事务支持}：MySQL提供完整的事务支持，确保数据的一致性和完整性。

        \item \textbf{成熟稳定}：MySQL是成熟的关系型数据库，具有丰富的功能和良好的性能表现。

        \item \textbf{TypeORM集成}：使用TypeORM作为ORM框架，支持TypeScript类型检查和数据库迁移。
    \end{itemize}

    \subsection{核心表结构}

    APPT平台共设计12个核心数据表，主要表结构如下：

    \begin{table}[htbp]
        \caption{APPT平台核心数据表}
        \label{tab:database-tables}
        \centering
        \begin{tabular}{lll}
            \toprule
            表名 & 中文名称 & 主要功能 \\
            \midrule
            users & 用户表 & 存储用户账户信息和认证数据 \\
            slide\_spaces & 幻灯片空间表 & 组织和管理幻灯片集合 \\
            slides & 幻灯片表 & 存储具体的幻灯片文档内容 \\
            slide\_versions & 版本控制表 & 记录幻灯片的版本历史 \\
            slide\_user\_roles & 用户角色表 & 定义用户在幻灯片中的权限 \\
            slide\_comments & 评论表 & 存储对幻灯片的评论和回复 \\
            snippets & 代码片段表 & 存储用户可复用的代码块 \\
            registration\_codes & 注册码表 & 管理用户注册码信息 \\
            slidev\_themes & Slidev主题表 & 存储Slidev主题配置 \\
            slidev\_plugins & Slidev插件表 & 存储Slidev插件配置 \\
            likes\_slide & 幻灯片点赞表 & 记录用户对幻灯片的点赞 \\
            likes\_space & 空间点赞表 & 记录用户对空间的点赞 \\
            \bottomrule
        \end{tabular}
    \end{table}

    \subsection{重点表结构详解}

    本节重点介绍 \texttt{slides}（幻灯片文档表）与 \texttt{slide\_comments}（评论表）两张表的详细字段设计。两表均以树形结构组织数据，但面对截然不同的使用场景，采取了差异显著的存储策略，具体对比详见本节末尾。

    \subsubsection{slides 幻灯片文档表}

    幻灯片文档表是平台的核心内容载体，存储每一个幻灯片文档的标题、正文与状态信息。\texttt{slides} 表通过 \texttt{parent\_id} 字段自关联，使空间内的文档可以形成任意深度的父子层级结构，支持知识库的树形目录组织，如表 \ref{tab:slides} 所示。

    \begin{table}[htbp]
        \caption{slides 幻灯片文档表结构}
        \label{tab:slides}
        \centering
        \begin{tabular}{llll}
            \toprule
            字段名 & 数据类型 & 约束 & 说明 \\
            \midrule
            id & bigint & PK，自增 & 文档唯一主键 \\
            title & varchar(255) & NOT NULL & 文档标题 \\
            content & longtext & NULLABLE & 幻灯片 Markdown 正文 \\
            slide\_space\_id & bigint & FK $\rightarrow$ slide\_spaces.id & 所属空间ID \\
            parent\_id & bigint & FK $\rightarrow$ slides.id & 父文档ID，根节点为NULL \\
            allow\_comment & tinyint(1) & DEFAULT 1 & 是否允许评论 \\
            created\_by & bigint & FK $\rightarrow$ users.id & 创建者用户ID \\
            is\_public & tinyint(1) & DEFAULT 0 & 是否公开（0私有，1公开）\\
            is\_build & tinyint(1) & DEFAULT 0 & 是否已编译生成预览 \\
            preview\_url & varchar(255) & NULLABLE & 编译后预览地址 \\
            created\_at & datetime & NOT NULL & 创建时间 \\
            updated\_at & datetime & NOT NULL & 最近修改时间 \\
            \bottomrule
        \end{tabular}
    \end{table}

    \texttt{parent\_id} 是文档树的核心字段，自关联至同表的 \texttt{id}，构成邻接表（Adjacency List）模型；根文档的 \texttt{parent\_id} 为 \texttt{NULL}。平台支持用户通过拖拽方式调整文档在目录树中的位置，即节点的父子关系可动态变更，因此 \texttt{slides} 表仅使用邻接表而不引入物化路径。\texttt{allow\_comment} 与 \texttt{is\_public} 字段提供文档级的细粒度权限控制，\texttt{is\_build} 与 \texttt{preview\_url} 则记录文档的编译构建状态，供前端判断是否展示实时预览或静态构建预览。

    \subsubsection{slide\_comments 评论表}

    评论表支持用户对幻灯片内容进行多级评论与回复，其树形结构通过\textbf{邻接表}与\textbf{物化路径}两种模式的组合实现，如表 \ref{tab:slide-comments} 所示。

    \begin{table}[htbp]
        \caption{slide\_comments 评论表结构}
        \label{tab:slide-comments}
        \centering
        \begin{tabular}{llll}
            \toprule
            字段名 & 数据类型 & 约束 & 说明 \\
            \midrule
            id & bigint & PK，自增 & 评论唯一主键 \\
            slide\_id & bigint & FK $\rightarrow$ slides.id & 所属幻灯片ID \\
            user\_id & bigint & FK $\rightarrow$ users.id & 评论发表者ID \\
            content & text & NOT NULL & 评论正文内容 \\
            path & varchar(255) & NULLABLE & 根节点到本节点的物化路径 \\
            reply\_id & bigint & FK $\rightarrow$ slide\_comments.id & 父评论ID，顶层评论为NULL \\
            created\_at & datetime & NOT NULL & 评论发布时间 \\
            \bottomrule
        \end{tabular}
    \end{table}

    评论表的关键设计要点如下：

    \begin{itemize}
        \item \textbf{邻接表（reply\_id）}：\texttt{reply\_id} 字段指向同表中父评论的 \texttt{id}，构成邻接表（Adjacency List）模型。顶层评论的 \texttt{reply\_id} 为 \texttt{NULL}，子评论的 \texttt{reply\_id} 指向其直接父节点。邻接表写入成本低，插入新评论只需设置父节点引用，无需修改其他记录，适合高频评论写入场景。

        \item \textbf{物化路径（path）}：\texttt{path} 字段存储从根节点到当前节点的完整路径字符串，格式如 \texttt{/1/5/12/}，其中每段数字为对应层级评论的 \texttt{id}。利用 SQL 前缀查询 \texttt{WHERE path LIKE '/1/5/\%'}，可在不执行递归查询的情况下一次性获取任意评论节点下的全部子孙评论；删除某条评论及其所有回复时同样只需一条带前缀匹配的 \texttt{DELETE} 语句，避免了递归删除带来的多次数据库往返。

        \item \textbf{权限联动控制}：评论功能受双重权限约束。\texttt{slides} 表中的 \texttt{allow\_comment} 字段由幻灯片所有者控制，决定该文档是否开放评论；用户角色须为 \texttt{commenter} 及以上级别（\texttt{commenter}、\texttt{editor}、\texttt{owner}）方可发表评论，\texttt{viewer} 角色仅有只读权限。
    \end{itemize}

    \subsubsection{两表设计思路对比}

    \texttt{slides} 与 \texttt{slide\_comments} 均以树形结构组织数据，但两者在节点移动、子树展开和删除三个维度上的使用模式截然不同，由此导致了存储策略的根本差异，具体如表 \ref{tab:design-contrast} 所示。

    \begin{table}[htbp]
        \caption{两表树形存储策略对比}
        \label{tab:design-contrast}
        \centering
        \begin{tabular}{lll}
            \toprule
            操作维度 & slides（文档树） & slide\_comments（评论树）\\
            \midrule
            节点结构调整 & 需要（支持拖拽排序） & 不需要（评论发出后不移动）\\
            子树展开方式 & 逐级展开（按需加载） & 一次性展开整个讨论串 \\
            子树删除方式 & 递归逐节点删除 & 单条 \texttt{LIKE} 前缀匹配删除 \\
            树形存储模式 & 仅邻接表（\texttt{parent\_id}）& 邻接表 + 物化路径双模式 \\
            \bottomrule
        \end{tabular}
    \end{table}

    \textbf{为何文档树不使用物化路径？} \texttt{slides} 表支持用户通过拖拽调整文档节点的位置，即节点的父子关系会动态变更。若引入物化路径，每次移动一个节点，都需要将该节点及其所有子孙节点的 \texttt{path} 字段全部更新，更新量与子树规模成正比，写放大严重。因此文档树仅维护 \texttt{parent\_id} 邻接表，移动节点只需修改一条记录的父节点引用，代价恒定。子树展开采用逐级按需加载策略，每次展开只查询当前层的直接子节点（\texttt{WHERE parent\_id = ?}），避免一次性加载过深的文档树；删除文档节点时则通过递归遍历依次删除子节点，与文档树规模相对较小、删除操作频率较低的实际使用场景相匹配。

    \textbf{为何评论树引入物化路径？} 评论一经发布，其所属的父评论不会改变——用户无法将一条回复"移动"到另一个讨论串下，因此物化路径的维护成本几乎为零，写入时仅需将父节点的 \texttt{path} 追加当前评论 \texttt{id} 即可生成自身路径。与此同时，评论的两类高频操作恰好是物化路径最擅长的场景：展开某条评论下的全部回复，只需 \texttt{WHERE path LIKE '/1/5/\%'} 一条查询即可一次性取出整个子树，无需递归；删除某条评论及其全部子孙回复，同样通过 \texttt{DELETE WHERE path LIKE '/1/5/\%'} 单条语句完成，避免了逐节点递归删除带来的多次数据库往返。物化路径以冗余一个 \texttt{path} 字段为代价，将子树查询和删除的复杂度从递归的 $O(n)$ 次查询降低为常数次，在评论这一结构固定、子树操作高频的场景下性价比极高。

    \subsection{核心实体关系图}

    APPT平台的核心实体关系如图 \ref{fig:er-diagram} 所示，展示了各表之间的关系：

    \clearpage
    \begin{figure}[!htbp]
        \centering
        \begin{tikzpicture}[
            entity/.style={
                rectangle, draw, rounded corners, fill=blue!10,
                align=center, minimum width=3.8cm, minimum height=1.8cm,
                text width=3.4cm, inner sep=0.4cm, font=\small
            },
            relationship/.style={
                diamond, draw, fill=green!10, align=center,
                minimum width=2.5cm, minimum height=1.4cm, inner sep=0.25cm, font=\small
            },
            arrow/.style={-, thick},
            label/.style={font=\small},
            scale=0.4, transform shape
        ]

        % ==================== 实体节点（松散布局 + 上移调整） ====================
        % 上层
        \node[entity] (user)   at (0, 15) {用户表\\users};
        \node[entity, right=10cm of user] (space) {幻灯片空间表\\slide\_spaces};

        % 中层（代码片段表上移）
        \node[entity] (snippet) at (-6.5, 8.5) {代码片段表\\snippets};
        \node[entity] (slide)   at (5, 7)     {幻灯片表\\slides};
        \node[entity, right=8.5cm of slide] (comment) {评论表\\slide\_comments};

        % 下层（注册码表大幅上移）
        \node[entity] (version) at (10, -1)    {版本控制表\\slide\_versions};
        \node[entity] (role)    at (-4, -1)    {用户角色表\\slide\_user\_roles};
        \node[entity] (regcode) at (-16.8, 8.8){注册码表\\registration\_codes};

        % ==================== 关系节点（菱形） ====================
        \node[relationship] (r_user_space)    at ($(user)!0.5!(space) + (0,-1.2)$) {拥有};
        \node[relationship] (r_space_slide)   at ($(space)!0.5!(slide) + (0,-1.2)$) {包含};
        \node[relationship] (r_slide_version) at ($(slide)!0.5!(version) + (0,-1.1)$) {版本化};
        \node[relationship] (r_user_slide)    at ($(user)!0.5!(slide) + (-2.5,-2)$) {协作};
        \node[relationship] (r_slide_comment) at ($(slide)!0.5!(comment) + (1.3,0)$) {评论};
        \node[relationship] (r_user_role)     at ($(user)!0.5!(role) + (-1.4,-1.7)$) {授权};
        \node[relationship] (r_user_regcode)  at ($(user)!0.5!(regcode) + (-2.2,-2.8)$) {使用};
        \node[relationship] (r_user_snippet)  at ($(user)!0.5!(snippet) + (0,-1.2)$) {创建};

        % ==================== 连线（全部改为直线，无箭头） ====================
        % 用户 1:N 拥有 空间
        \draw[arrow] (user) -- node[pos=0.35, above, label] {1} (r_user_space);
        \draw[arrow] (r_user_space) -- node[pos=0.65, above, label] {N} (space);

        % 空间 1:N 包含 幻灯片
        \draw[arrow] (space) -- node[pos=0.35, right, label] {1} (r_space_slide);
        \draw[arrow] (r_space_slide) -- node[pos=0.65, right, label] {N} (slide);

        % 幻灯片 1:N 版本化
        \draw[arrow] (slide) -- node[pos=0.35, right, label] {1} (r_slide_version);
        \draw[arrow] (r_slide_version) -- node[pos=0.65, right, label] {N} (version);

        % 用户 1:N 协作 幻灯片
        \draw[arrow] (user) -- node[pos=0.35, left, label] {1} (r_user_slide);
        \draw[arrow] (r_user_slide) -- node[pos=0.65, above, label] {N} (slide);

        % 幻灯片 1:N 评论
        \draw[arrow] (slide) -- node[pos=0.4, above, label] {1} (r_slide_comment);
        \draw[arrow] (r_slide_comment) -- node[pos=0.6, above, label] {N} (comment);

        % 用户 1:N 授权 角色
        \draw[arrow] (user) -- node[pos=0.35, left, label] {1} (r_user_role);
        \draw[arrow] (r_user_role) -- node[pos=0.65, right, label] {N} (role);

        % 用户 1:N 创建 代码片段
        \draw[arrow] (user) -- node[pos=0.4, left=0.4cm, label] {1} (r_user_snippet);
        \draw[arrow] (r_user_snippet) -- node[pos=0.6, right, label] {N} (snippet);

        % 用户 1:N 使用 注册码（直线连接）
        \draw[arrow] (user) -- node[pos=0.38, left=0.5cm, label] {1} (r_user_regcode);
        \draw[arrow] (r_user_regcode) -- node[pos=0.65, left, label] {N} (regcode);

        \end{tikzpicture}
        \caption{APPT平台核心实体关系图}
        \label{fig:er-diagram}
    \end{figure}
    主要关系包括：
    \begin{itemize}
        \item \textbf{用户与空间}：一个用户可以拥有多个幻灯片空间（一对多关系）。
        \item \textbf{空间与幻灯片}：一个空间可以包含多个幻灯片（一对多关系）。
        \item \textbf{幻灯片与版本}：一个幻灯片可以有多个版本（一对多关系）。
        \item \textbf{用户与权限}：一个用户可以在多个幻灯片中拥有不同的权限（多对多关系）。
        \item \textbf{幻灯片与评论}：一个幻灯片可以有多个评论（一对多关系）。
        \item \textbf{评论与回复}：评论之间可以形成回复关系（自关联一对多关系）。
    \end{itemize}

    \section{核心功能实现}\label{sec:core-features}

    \subsection{实时协作编辑功能}

    实时协作编辑功能基于Yjs CRDT算法实现，具体实现细节如下：

    \begin{itemize}
        \item \textbf{文档结构设计}：将幻灯片文档设计为Yjs的YText类型，支持富文本编辑和格式标记。

        \item \textbf{操作同步机制}：通过WebSocket连接同步编辑操作，使用增量更新减少网络传输量。

        \item \textbf{冲突解决策略}：采用CRDT算法自动解决并发编辑冲突，确保最终一致性。

        \item \textbf{性能优化}：使用操作批处理和压缩算法优化网络传输性能，提高响应速度。
    \end{itemize}

    \subsection{协同算法原理}

    APPT平台采用基于CRDT（Conflict-Free Replicated Data Type，无冲突复制数据类型）的协同算法实现实时协作编辑。CRDT是一种分布式数据结构的数学理论，能够确保在异步网络环境下，多个副本之间的数据最终保持一致，无需中央协调器解决冲突。

    \subsubsection{CRDT基本概念}

    CRDT算法主要分为两种类型：
    \begin{itemize}
        \item \textbf{基于操作的CRDT（Op-based CRDT）}：通过同步操作序列实现一致性，每个操作都是可交换、可结合和幂等的。
        \item \textbf{基于状态的CRDT（State-based CRDT）}：通过同步完整状态或状态差异实现一致性，状态合并操作满足交换律、结合律和幂等律。
    \end{itemize}

    APPT平台采用Yjs库实现的基于操作的CRDT算法，具体特点如下：
    \begin{itemize}
        \item \textbf{交换性}：操作顺序不影响最终结果，无论操作到达顺序如何，最终状态一致。
        \item \textbf{结合性}：操作组合方式不影响最终结果，支持分组传输和批处理优化。
        \item \textbf{幂等性}：重复执行同一操作不会改变状态，确保网络重传和重复消息不会导致数据不一致。
    \end{itemize}

    \subsubsection{Yjs文档模型设计}

    Yjs采用分层文档模型设计，核心数据结构包括：
    \begin{itemize}
        \item \textbf{YText}：富文本数据类型，支持字符级别的插入、删除和格式属性设置，用于表示幻灯片内容。
        \item \textbf{YArray}：有序数组类型，支持元素的插入、删除和移动，用于管理幻灯片页面顺序。
        \item \textbf{YMap}：键值对类型，支持属性的增删改查，用于存储幻灯片元数据和用户配置。
        \item \textbf{Transaction}：事务机制，将多个操作打包为原子操作单元，确保操作的完整性。
    \end{itemize}

    文档状态通过逻辑时钟（Lamport Timestamp）进行版本控制，每个操作附带唯一的操作ID（客户端ID + 逻辑时钟），确保全局有序性。

    \subsubsection{冲突解决机制}

    当多个用户同时编辑同一文档位置时，CRDT算法通过以下机制自动解决冲突：
    \begin{enumerate}
        \item \textbf{位置标识算法}：使用唯一标识符（ID）标记每个字符位置，而不是传统的位置索引。字符插入时生成新的唯一ID，删除时标记为墓碑状态。
        \item \textbf{偏序关系定义}：基于操作ID的因果关系定义偏序关系，确保因果一致性。如果操作A在操作B之前发生，且操作B依赖于操作A的结果，则A必须优先于B执行。
        \item \textbf{合并策略}：当并发操作无因果关系时，采用确定性合并策略。例如，字符插入冲突时，基于操作ID的字典序确定最终位置顺序。
    \end{enumerate}

    以文本编辑为例，假设用户A在位置P插入字符"X"，用户B在相同位置插入字符"Y"。传统方法会产生冲突，而CRDT算法通过以下步骤解决：
    \begin{enumerate}
        \item 为"X"生成唯一ID：\texttt{(clientA, timestamp1)}
        \item 为"Y"生成唯一ID：\texttt{(clientB, timestamp2)}
        \item 基于ID排序确定最终顺序：若\texttt{clientA < clientB}，则最终文本为"XY"；否则为"YX"
    \end{enumerate}

    \subsubsection{网络同步协议}

    APPT平台采用优化的网络同步协议提高协作效率：
    \begin{itemize}
        \item \textbf{增量同步}：仅传输文档状态差异，而非完整文档内容，减少网络带宽消耗。
        \item \textbf{操作压缩}：将连续的编辑操作合并为批量更新，减少网络往返次数。
        \item \textbf{状态向量（State Vector）}：记录每个客户端已知的最新状态，用于计算最小差异集。
        \item \textbf{断线重连恢复}：通过状态向量对比，快速同步断线期间错过的更新。
    \end{itemize}

    \subsubsection{性能优化策略}

    为提高大规模文档的协作性能，APPT平台实施以下优化：
    \begin{itemize}
        \item \textbf{分层更新}：将文档划分为逻辑块（如幻灯片页面），仅同步当前编辑块的变更。
        \item \textbf{延迟同步}：对高频率编辑操作（如连续输入）进行去抖动处理，减少同步频率。
        \item \textbf{内存管理}：定期清理历史操作记录，释放内存资源，同时保留必要的版本信息。
        \item \textbf{本地优先}：编辑操作先在本地应用并确认，再异步同步到服务器，提供即时反馈体验。
    \end{itemize}

    \subsubsection{与传统协同算法的对比}

    与传统基于操作转换（OT）的协同算法相比，CRDT算法具有以下优势：
    \begin{table}[htbp]
        \caption{CRDT与OT算法对比}
        \label{tab:crdt-vs-ot}
        \centering
        \begin{tabular}{lll}
            \toprule
            对比维度 & CRDT算法 & OT算法 \\
            \midrule
            冲突解决 & 自动解决，无需中央协调器 & 需要中央服务器协调转换 \\
            网络要求 & 支持P2P直接同步 & 必须通过中央服务器 \\
            算法复杂度 & 相对简单，基于数学理论 & 复杂，需要维护转换函数 \\
            扩展性 & 容易扩展到大规模节点 & 中央服务器可能成为瓶颈 \\
            离线支持 & 天然支持，本地操作独立 & 需要复杂的状态同步机制 \\
            \bottomrule
        \end{tabular}
    \end{table}

    CRDT算法的数学基础确保了系统的正确性和可靠性，使得APPT平台能够在分布式环境下提供稳定、高效的实时协作体验。

    \subsubsection{文档协同流程图}

    APPT平台的文档协同流程采用条件分支设计，根据文档的协作状态动态选择独立编辑模式或实时协作模式。具体流程如图 \ref{fig:collaboration-flowchart} 所示。

    \begin{figure}[htbp]
        \centering
        \begin{tikzpicture}[
            node distance=1.0cm and 1.8cm,
            startstop/.style={rectangle, rounded corners, draw, fill=red!10, align=center, minimum width=2.2cm, minimum height=0.9cm, font=\scriptsize},
            process/.style={rectangle, draw, fill=blue!10, align=center, minimum width=2.2cm, minimum height=0.9cm, font=\scriptsize},
            decision/.style={rectangle, draw, fill=green!10, align=center, minimum width=2.2cm, minimum height=1.0cm, font=\scriptsize, text width=1.8cm},
            io/.style={rectangle, draw, fill=yellow!10, align=center, minimum width=2.2cm, minimum height=0.9cm, font=\scriptsize},
            connector/.style={circle, draw, fill=orange!10, align=center, minimum width=0.7cm, minimum height=0.7cm, font=\scriptsize},
            arrow/.style={->, >=stealth, thick},
            label/.style={font=\scriptsize},
            scale=0.75
        ]

        % 定义节点
        \node[startstop] (start) {用户打开文档};
        \node[io, below=of start] (load) {加载文档内容\\与权限信息};
        \node[decision, below=of load, yshift=-0.3cm] (check) {文档有\\协作者？};

        % 左侧分支：无协作者（独立模式）
        \node[process, left=of check, xshift=-1.2cm] (localEdit) {直接本地编辑\\使用CodeMirror};
        \node[process, below=of localEdit, yshift=-1.2cm] (localSave) {定期保存到\\后端数据库};
        \node[connector, below=of localSave, yshift=-1.2cm] (localEnd) {结束};

        % 右侧分支：有协作者（协同模式）
        \node[process, right=of check, xshift=1.2cm] (initWs) {初始化WebSocket\\连接Hocuspocus};
        \node[process, below=of initWs, yshift=-0.5cm] (initYjs) {创建Yjs文档实例\\绑定CodeMirror};
        \node[process, below=of initYjs, yshift=-0.5cm] (sync) {同步文档状态\\使用CRDT算法};
        \node[process, below=of sync, yshift=-0.5cm] (edit) {协同编辑\\实时冲突解决};
        \node[process, below=of edit, yshift=-0.5cm] (persist) {保存到\\后端数据库};
        \node[connector, below=of persist, yshift=-0.5cm] (collabEnd) {结束};

        % 连接线
        \draw[arrow] (start) -- (load);
        \draw[arrow] (load) -- (check);

        % 左侧分支连接
        \draw[arrow] (check) -- node[above, label] {否} (localEdit);
        \draw[arrow] (localEdit) -- (localSave);
        \draw[arrow] (localSave) -- (localEnd);

        % 右侧分支连接
        \draw[arrow] (check) -- node[above, label] {是} (initWs);
        \draw[arrow] (initWs) -- (initYjs);
        \draw[arrow] (initYjs) -- (sync);
        \draw[arrow] (sync) -- (edit);
        \draw[arrow] (edit) -- (persist);
        \draw[arrow] (persist) -- (collabEnd);

        % 添加虚线框区分模式
        % 独立编辑模式虚线框
        \coordinate (localBoxNW) at ($(localEdit.north west)+(-0.3,0.3)$);
        \coordinate (localBoxSE) at ($(localSave.south east |- localEnd.south)+(0.3,-0.3)$);
        \draw[dashed, gray, thick] (localBoxNW) rectangle (localBoxSE);
        \coordinate (localLabelPos) at ($(localBoxNW)!0.5!(localBoxSE |- localBoxNW)$);
        \node[gray, above=0.2cm] at (localLabelPos) {\textbf{独立编辑模式}};

        % 实时协作模式虚线框
        \coordinate (collabBoxNW) at ($(initWs.north west)+(-0.3,0.3)$);
        \coordinate (collabBoxSE) at ($(persist.south east |- collabEnd.south)+(0.3,-0.3)$);
        \draw[dashed, gray, thick] (collabBoxNW) rectangle (collabBoxSE);
        \coordinate (collabLabelPos) at ($(collabBoxNW)!0.5!(collabBoxSE |- collabBoxNW)$);
        \node[gray, above=0.2cm] at (collabLabelPos) {\textbf{实时协作模式}};

        \end{tikzpicture}
        \caption{APPT平台文档协同流程图}
        \label{fig:collaboration-flowchart}
    \end{figure}

    流程的关键决策点在于文档是否有协作者，具体实现逻辑如下：

    \begin{enumerate}
        \item \textbf{文档加载阶段}：用户打开文档时，系统加载文档内容和权限信息，同时查询该文档的协作者列表。

        \item \textbf{协作状态判断}：检查文档是否有除当前用户外的其他协作者。如果文档为私有文档或只有当前用户有访问权限，则进入独立编辑模式；否则进入实时协作模式。

        \item \textbf{独立编辑模式}：适用于个人文档或私有编辑场景：
        \begin{itemize}
            \item 直接使用CodeMirror编辑器进行本地编辑，无网络同步开销。
            \item 编辑操作实时应用，提供流畅的本地编辑体验。
            \item 定期或手动保存时将内容同步到后端MySQL数据库。
        \end{itemize}

        \item \textbf{实时协作模式}：适用于多人协作场景：
        \begin{itemize}
            \item 初始化WebSocket连接，连接到Hocuspocus协作服务器。
            \item 创建Yjs文档实例，并与CodeMirror编辑器绑定。
            \item 通过CRDT算法自动同步文档状态，解决编辑冲突。
            \item 支持实时光标追踪、用户状态显示等协作功能。
            \item 定期将协作文档状态持久化到后端数据库，防止数据丢失。
        \end{itemize}

        \item \textbf{模式切换机制}：当文档协作状态发生变化时（如新增协作者或协作者全部离开），系统可动态切换编辑模式，确保用户体验的一致性。
    \end{enumerate}

    这种条件分支设计具有以下优势：
    \begin{itemize}
        \item \textbf{资源优化}：私有文档无需建立WebSocket连接，节省服务器资源。
        \item \textbf{性能提升}：独立编辑模式减少网络延迟，提供更快的响应速度。
        \item \textbf{灵活适应}：根据实际协作需求动态调整技术方案。
        \item \textbf{用户体验}：保持编辑界面一致性，用户无需关心底层技术差异。
    \end{itemize}

    APPT平台通过智能的协作状态判断和模式选择机制，在保证功能完整性的同时，优化系统性能和用户体验，实现了高效、灵活的文档协同编辑解决方案。

    \subsection{权限管理系统}

    权限管理系统基于RBAC模型实现，具体实现细节如下：

    \begin{itemize}
        \item \textbf{角色定义}：定义Owner、Editor、Commenter、Viewer四种角色，每种角色具有不同的权限集合。

        \item \textbf{权限分配}：通过slide\_user\_roles表记录用户在特定幻灯片中的角色和权限。

        \item \textbf{权限验证}：在API接口和WebSocket连接中验证用户权限，阻止未授权操作。

        \item \textbf{文档可见性}：支持公开和私有两种文档可见性设置，控制文档的访问范围。
    \end{itemize}

    \subsection{AI辅助功能}

    AI辅助功能集成大语言模型API，具体实现细节如下：

    \begin{itemize}
        \item \textbf{API集成}：集成通义千问等大语言模型API，支持文本生成和图像生成。

        \item \textbf{提示词工程}：设计专门的提示词模板，提高AI生成内容与Slidev语法的适配性。

        \item \textbf{上下文管理}：维护对话上下文，支持多轮交互和上下文相关的建议生成。

        \item \textbf{结果处理}：对AI生成的内容进行格式化和语法检查，确保符合Slidev规范。
    \end{itemize}

    \subsection{版本控制系统}

    版本控制系统设计简化操作流程，具体实现细节如下：

    \begin{itemize}
        \item \textbf{自动保存}：定期自动保存文档内容，防止数据丢失。

        \item \textbf{手动版本创建}：支持用户手动创建重要版本，并添加提交消息说明。

        \item \textbf{差异对比}：使用diff算法计算版本间差异，并以可视化方式展示变化内容。

        \item \textbf{一键回滚}：支持将文档内容快速回滚到任意历史版本。
    \end{itemize}

    \section{测试与评估}\label{sec:testing-evaluation}

    \subsection{测试方法}

    为确保系统质量，采用以下测试方法：

    \begin{itemize}
        \item \textbf{单元测试}：对核心功能模块进行单元测试，确保代码逻辑正确性。

        \item \textbf{集成测试}：测试模块间的集成和交互，确保系统整体功能正常。

        \item \textbf{用户体验测试}：邀请真实用户使用系统，收集反馈意见并改进用户体验。
    \end{itemize}

    \subsection{测试结果}

    测试结果表明，APPT平台在功能、性能和用户体验方面表现良好：

    \begin{itemize}
        \item \textbf{功能测试}：所有核心功能均能正常工作，实时协作、权限管理、AI辅助等功能达到设计要求。

        \item \textbf{兼容性测试}：系统在现代主流浏览器（Chrome、Firefox、Safari、Edge）上均能正常运行。

        \item \textbf{用户体验测试}：用户反馈表明，系统界面友好，操作流程清晰，学习成本较低。
    \end{itemize}

    \subsection{性能评估}

    对系统关键性能指标进行评估，结果如下：
% TODO添加协作编辑性能评估数据
    \begin{table}[htbp]
        \caption{APPT平台性能评估结果}
        \label{tab:performance-evaluation}
        \centering
        \begin{tabular}{lll}
            \toprule
            性能指标 & 测试条件 & 结果 \\
            \midrule
            页面加载时间 & 首次访问 & 2.1秒 \\
            编辑器响应时间 & 普通编辑操作 & 120毫秒 \\
            AI生成响应时间 & 中等长度请求 & 3.5秒 \\
            \bottomrule
        \end{tabular}
    \end{table}

    \section{小结}\label{sec:design-summary}

    本章详细介绍了APPT平台的设计目标、系统架构、前后端设计、数据库设计以及核心功能实现。通过前后端分离的架构、基于CRDT的实时协作技术、RBAC权限管理系统和AI辅助功能集成，APPT平台实现了智能幻灯片协作编辑的核心功能。

    系统设计充分考虑了用户体验和技术实现的平衡，在保证功能完整性的同时，注重降低非技术用户的使用门槛。测试结果表明，系统在功能、性能和用户体验方面均达到设计要求，能够满足技术团队和教育工作者对智能幻灯片协作平台的需求。

    未来，我们将继续优化系统性能，扩展AI辅助功能，增加更多协作场景支持，并根据用户反馈持续改进系统功能和用户体验，使APPT平台成为更加强大和易用的幻灯片协作工具。

\end{cuzchapter}